\chapter{Intraosseous Access}

\section*{Overview}
Intraosseous (IO) access is the placement of a needle directly into the medullary cavity of a bone to provide rapid vascular access in emergencies. It allows infusion of fluids and medications when intravenous cannulation fails or is impractical.

\section*{Alternative Names}
IO line, intraosseous infusion, intraosseous cannulation, bone marrow needle placement

\section*{Principle}
The highly vascular bone marrow provides a non-collapsible venous plexus that communicates directly with the systemic circulation. Infused fluids and drugs reach the central circulation reliably even in severe hypovolemia or cardiac arrest, when peripheral veins are collapsed.

\section*{History}
Intraosseous infusion was described as early as the 1920s and was used widely during World War II and in pediatric medicine. Its use has expanded to adult emergency care, trauma, and prehospital settings with modern powered insertion devices.

\section*{Why It Is Done}
Performed when rapid vascular access is needed and peripheral or central intravenous access fails, such as in cardiac arrest, severe shock, trauma, burns, status epilepticus, or sepsis in infants and adults.

\section*{Is This a Primary or Secondary Test or Procedure}
A secondary rescue procedure used when standard intravenous access is unavailable, but a primary approach in pediatric resuscitation protocols.

\section*{How It Is Performed}
The site, usually the proximal tibia, proximal humerus, or distal femur, is identified and prepared with antiseptic. A manual or powered needle is inserted through the cortex into the marrow cavity with a popping sensation, the stylet is removed, and aspiration of blood or marrow confirms placement. The line is secured and flushed, then connected to fluids and medications.

\section*{Preparation}
Informed consent when feasible, and preparation of the insertion site with sterile antiseptic. No fasting or laboratory testing is required in the emergency setting.

\section*{Contraindications and Precautions}
Contraindications include fracture of the target bone, recent IO placement in the same bone, infection or burns at the site, and severe osteoporosis in some cases. Precautions include avoiding extravasation of medications, confirming placement by aspiration, and removing the line within 24 hours or once venous access is established.

\section*{Risks}
Extravasation and compartment syndrome, osteomyelitis, fracture, fat embolism, and rarely injury to neurovascular structures.

\section*{Aftercare and Recovery}
The IO line is used for resuscitation while standard intravenous or central venous access is obtained. It is removed as soon as alternative access is available, and the site is covered with a sterile dressing and observed for infection.

\section*{Outcome and Follow-up}
The procedure is successful when aspiration or flushing confirms placement and fluids and medications infuse freely without significant resistance or extravasation. Failure to infuse suggests malposition, occlusion, or dislodgement requiring reinsertion.

\section*{Expected Results}
Free aspiration of blood or marrow, a stable insertion site without swelling, and successful delivery of fluids and medications constitute a satisfactory result.

\section*{Follow-up}
The patient is transitioned to intravenous access and monitored for local complications such as extravasation, infection, or compartment syndrome. Radiographic or laboratory confirmation of position may be obtained if malposition is suspected.

\section*{References}
\begin{enumerate}
  \item MedlinePlus. Intraosseous infusions. U.S. National Library of Medicine; 2025.
  \item Current Medical Diagnosis and Treatment. McGraw-Hill; 2025.
\end{enumerate}

\clearpage
